\begin{tikzpicture}[font=\scriptsize, >=latex]
% process address space, low->high
\draw[thick] (0.6,0.4) rectangle (2.4,4.8);
\foreach \y/\t/\c in {
  4.4/{high: stack (grows down)}/{graybg},
  3.5/{mmap region}/{white},
  2.7/{heap (grows up)}/{graybg},
  2.1/{.bss + .data}/{white},
  1.5/{.text}/{accentsoft},
  0.8/{low: unused + kernel-reserved}/{white}}{
  \draw[fill=\c] (0.6,\y-0.35) rectangle (2.4,\y+0.35);
  \node[font=\tiny] at (1.5,\y) {\t};
}
\node[right, align=left, font=\tiny, text width=6.8cm] at (3.0,3.5)
{the ELF's PT\_LOAD segments become text+data here};
\node[right, align=left, font=\tiny, text width=6.8cm] at (3.0,2.7)
{the loader builds this space; the CPU only sees the addresses};
\node[right, align=left, font=\tiny, text width=6.8cm] at (3.0,1.5)
{entry point jumps here --- your cpu8's pc=0x0, grown up};
\node[align=left, text=slate, text width=10cm] at (5.3,-0.4)
{Same layout idea as the machine you built in Book 5, same reason:
the hardware reads addresses, the software draws the map.};
\end{tikzpicture}
